\section{Conclusion}

This thesis designed, implemented, and evaluated an adaptive CCTV compression system using dynamic ROI control. RQ1 is answered affirmatively as a construction result: a deployable YOLOv8n/FFmpeg edge-cloud prototype with heuristic risk scoring (NORMAL/ALERT/CRITICAL), Lagrangian-inspired ROI allocation, 15-second pre-event and 10-second post-event buffering with SHA-256 hash-chained manifests, and a Next.js monitoring dashboard was built and documented with all evaluated parameters. RQ2 is answered descriptively: the risk score tracks detected-object/ROI importance across CDnet-2014 and UCF-Crime probes (e.g.\ highway risk $\approx$ 0.563, pedestrians $\approx$ 0.492, night $\approx$ 0.024; Normal\_050 $\approx$ 0.994 versus Fighting002 $\approx$ 0.068), while group-aware calibration over 64 labelled seconds shows in-sample separability (descriptive AUC $\approx$ 0.9962, F1 $\approx$ 0.919) without established generalization. RQ3 is answered negatively: under the locked same-codec protocol, the current ROI construction did not outperform uniform encoding at matched bitrate (proposed ROI-PSNR $\approx$ 31.5 dB versus uniform $\approx$ 34.7 dB at $\approx$ 710--840 kbps), with BD-rate withheld. RQ4 is answered positively within scope: the event buffer demonstrated exact-byte tamper-evident integrity behaviour (5/5 direct tampers detected, zero clean-case false alarms) without any legal-admissibility claim. The overall contribution is therefore a deployable YOLOv8n/FFmpeg edge-cloud CCTV prototype with heuristic risk scoring, Lagrangian-inspired ROI allocation, and tamper-evident event buffering, evaluated through public-dataset probing, group-aware calibration analysis, controlled forensic tests, and a locked same-codec baseline comparison---with the current ROI construction's limitation stated as plainly as its capabilities.

\section{Limitations}

The work is limited by: a small calibration set with positives from a single video; single-class LOVO folds and one degenerate fold; hour-fixing in Phase~1; a single primary motion clip in Phase~4; a limited tamper-operation set with exact-byte scope and no legal validation; CPU-only execution with no energy or FPS claims; no independent detector ground truth in Phase~4 (consistency, not accuracy); withheld BD-rate; the dropped-hysteresis deviation of Phase~4 from the live system; and no systematic server/dashboard performance evaluation. Superseded exploratory figures (e.g.\ -87.95\%/-83.15\% BD-rate, 20.4 FPS, 217 kbps, 0.913 SSIM, greater than 95\% recall as a headline) are demoted and not claimed.

\section{Future Work}

Future work should: (i) redesign the ROI construction for moving-camera content (motion-compensated or learned backgrounds \cite{Zhao2025ASVC,Wu2024ROIDVC}, adaptive dilation, per-region QP optimization against ROI-weighted distortion); (ii) extend evaluation to more scenes, codecs (including AV1/VP9 paths already in the harness), and longer sequences with overlapping rate--distortion coverage sufficient for valid BD-rate computation; (iii) collect multi-video calibration labels so that generalization---not just apparent fit---can be tested under LOVO; (iv) broaden forensic testing (more tamper types, robustness-oriented integrity such as perceptual hashing) while preserving the exact-byte baseline; (v) re-examine the risk/anomaly relationship with frame-level anomaly annotations rather than clip labels, informed by the anomaly literature \cite{Duong2023Survey,ACMTaxonomy2024,Baradaran2021Anomaly}; and (vi) evaluate server/dashboard performance and energy consumption with instrumented metering. Each direction builds directly on the prototype and protocols documented in this book.
